% !TeX root = ../papers/nc-p3.tex
\section{四周期二次PBW表示で可能な先頭語グラフの一意性}
\label{sec:periodic-pbw-graph}
各矢空間に順序付き四文字を取り、基底と順序は四位置ごとに繰り返す。
固定した始点・長さで左辞書順を使う。ここでのPBW表示は、窓の六関係を
この順序で掃き出したものが二次Gröbner基底になるという意味である。
\begin{theorem}\label{thm:periodic-pbw-graph}
このようなPBW表示を持つ標準14641型AS窓の二次先頭語代数は、
位置ごとの文字の付け替えを除き一種類である。
具体的には各位置の文字に $1,2,3,4$ を付けると、位置 $p$ の $i$ と
位置 $p+1$ の $j$ を連ねた二文字が正規であることと $i\le j$ が同値になる。
この番号は許容遷移のグラフから一意に回復できる。
これは文字の番号であり、最初の辞書順における文字の大小との一致を要しない。
\end{theorem}
\begin{proof}
頂点を（位置、文字）の十六組とし、正規な二文字に対応する有向辺だけを取る。
各位置から次の位置へ辺が十本ある。正規語はこのグラフの歩道に等しく、
各始点で長さ $n\ge1$ の正規語数は $\binom{n+3}{3}$ である。

閉路を含む強連結成分は単純な有向サイクルでなければならない。
そうでないとある頂点から成分内へ出る辺が少なくとも二本あり、
その二辺から始まり元の頂点へ戻る二つの閉歩道がある。
反復して同じ長さにすれば、これらの任意の連接が相異なる歩道を与え、指数増大に反する。
成分を縮約したグラフは非巡回である。一つの歩道が通るサイクル成分数の最大を $s$ とする。
通る成分列、成分間の辺、サイクルへの入口と出口（端点を含む）を固定する方法は有限である。
残る自由度は各サイクルを余分に何周するかだけで、指定した全長について最後の周回数は
他の周回数から高々一通りに決まる。従って $s\ge1$ のとき全歩道数は $O(n^{s-1})$ であり、
$s=0$ なら長い歩道はない。三次増大により一列に少なくとも四サイクルが必要である。

辺は必ず一位置だけ進むので各サイクルの長さは4の正の倍数である。
十六頂点しかないから、先ほどの列は長さ4の四サイクルで全頂点を使い切る。
この列の到達順に番号を付ける。各サイクルはどの位置にも一つの文字を持ち、
成分をまたぐ辺は小さい番号から大きい番号へのみ進む。
各位置で同番号の四辺以外に可能なのは六対だけである。
辺が十本あるから六対が全て現れる。成分間の到達順は全順序なので番号も一意である。
\end{proof}

\begin{corollary}
標準有限双対Frobenius条件を満たし、全始点の三次経路空間の次元が20の四窓を取る。
四周期の順序付き基底による二次PBW表示が存在することは、基底と各位置の番号付けを選んで、
六関係の先頭二文字をちょうど $i>j$ の六対にできることと同値である。
この枝の全入力は標準Koszul左右AS正則である。
\end{corollary}
\begin{proof}
二次PBWならKoszulであり、有限双対Frobenius条件と標準型からAS性と標準Hilbert級数を得る。
従って定理が必要条件を与える。逆に、番号が非減少な三文字は重複組合せにより
$\binom{6}{3}=20$ 個である。停止する二次簡約でこれらは三次商を張り、次元から独立である。
従って長さ3の全ての重なりが解消する。第\ref{sec:two-window-ext}節の
ダイヤモンド補題の議論で全次数のPBW性、従って正則性が従う。

基底の存在は有限の方程式と非消滅条件にできる。四つの基底変換の行列式を反転し、
各位置で辞書順の文字とサイクル番号の対応を四文字の順列から選ぶ。
関係の $6\times16$ 係数行列の列を辞書順の降順に並べると、最初の $k$ 列の階数は、
その $k$ 列に含まれる指定先頭対の数でなければならず、これで主元集合も決定する。
この有限個の小行列条件を基底変換について射影すればよい。
ここで決めたのは先頭語代数とこのPBW枝への所属である。
同じ先頭対を持つ窓の下位係数には区別すべき添字保持同型類が残り、
PBW表示を持たない窓も別途扱う必要がある。
\end{proof}

\begin{corollary}\label{cor:pbw-nonreduced-four-cycle}
標準14641型窓において四辺の双線型零点スキーム
$E_p\subset\mathbb P(V_p^*)\times\mathbb P(V_{p+1}^*)$ が全て有限であるとする。
四周期の左辞書順で二次PBW表示があれば、四つの点 $[u_p]\in\mathbb P(V_p^*)$
が存在し、$(u_p,u_{p+1})$ は各 $E_p$ の\emph{非被約点}である（$p$ は法4）。
特に有限四辺の少なくとも一辺が被約なら、そのようなPBW表示は存在しない。
\end{corollary}
\begin{proof}
$R_p\subset V_p\otimes V_{p+1}$ を六関係とし、$u_p$ を辞書順で最大の文字の双対とする。
$u_p$ における第一因子の縮約 $L:R_p\to V_{p+1}$ の階数は、最大文字から始まる
先頭対の数である。実際、降順に並べた係数行列の最初の四列の階数に等しい。
この文字のサイクル番号が $i$ なら定理\ref{thm:periodic-pbw-graph}より階数は $i-1$。
一方、階数が2以下なら $u_p$ の上の零点対応は少なくとも射影直線を含むので
有限性に反する。従って四位置の最大文字の番号は全て4で、$L$ の階数は3である。
両位置の最大文字を並べた対は正規で、しかも係数行列の先頭列に当たる。
この列は零でなければならず、$(u_p,u_{p+1})\in E_p$ を得る。

右縮約を $D:R_p\to V_p$, $r\mapsto(1\otimes u_{p+1})(r)$ と書く。
$\ker L$ の元は最大文字から始まる全係数が零である。このとき次に大きい列は
「二番目の文字、次位置の最大文字」の対である。この対も第二の番号が4なので正規であり、
$\ker L$ のいずれの行でもこの列は零となる。さもなければその行自体が正規対を先頭語に持つ。
従って $D(\ker L)$ は残る二文字の張る空間に入り、
$\dim\ker L=3$ から $\ker L\cap\ker D\ne0$ を得る。

この共通核の非零関係は、点 $(u_p,u_{p+1})$ で両因子方向への微分がともに零である。
六本の定義方程式のヤコビ行列の階数は高々5であり、滑らかな六次元の積空間内で
$E_p$ の接空間は少なくとも一次元となる。有限スキームの幾何点が被約ならその局所環は体で
接空間は零である。従って得られた四点は全て非被約である。
\end{proof}

例えば第\ref{sec:finite-point-chart}節の $v_i^2-v_j^2$ による対称AS窓は
被約な二十点を持つため、このPBW枝の外にある。
\path{computations/verify_finite_square_chart.py} は二十組の評価から六関係を回復し、
各関係行列 $R$ の微分 $((R y)^{\mathsf t},x^{\mathsf t} R)$ を用いて全二十点で階数6を直接確認する。

\begin{corollary}\label{cor:pbw-tangent-flags}
有限四辺の場合、候補四周期列を一つ固定し、各辺の点における接空間次元を $t_p$ とする。
辞書順で最初の二つの双対文字が張る旗の二平面の候補は、各位置で
$\mathbb P^{t_p-1}$ に絞られる。特に全て $t_p=1$ なら、有限個の候補点列ごとに
高々 $(\mathbb P^1)^4$ 上で有限個の行列階数条件を調べれば、この表示の存在が決まる。
\end{corollary}
\begin{proof}
前系の $u=u_p$, $v=u_{p+1}$, $L$, $D$ を使う。積の射影接空間での接方程式は
接ベクトルを $\dot u\bmod ku$, $\dot v\bmod kv$ で表して
\[
             \dot u(D r)+\dot v(L r)=0\qquad(r\in R_p)
\]
である。$L$ の像は三次元の $\ker v$ なので、$\dot v\bmod kv$ は $\dot u$ から
一意に決まり、その存在条件は $\dot u(D\ker L)=0$ である。従って接空間の第一射影は
$(D\ker L)^\perp/ku$ に同型である。前系の第二列に関する議論から、
二番目の双対文字はまさにこの消滅条件を満たす必要がある。
$u$ を含む二平面の候補はこの商の直線なので主張の射影空間を得る。

$t_p=1$ ならこの二平面は一意である。残る三平面はその二平面を含むもの全体、
すなわち二次元商の直線からなる。旗を固定すれば具体的な適合基底の選び直しは不要である。
降順に並べた双対基底の変更は各文字にその前の文字の線型結合を加える操作と
非零スカラー倍で尽きる。次位置とテンソル積を取った辞書順の列でも
各先頭区間の張る空間が保たれるので、先の小行列階数は旗だけで決まる。
最大文字の番号は4に固定し、残る番号の対応を各位置で六通りから選べば十分である。
有限辺には有限個の幾何点しかないから、四周期列も有限集合である。
この判定には非被約点の接線を持つことの十分性は仮定しない。
\end{proof}

特に $t_p=t_{p+1}=1$ で、隣接する二位置の第二双対文字の候補を
$u'_p,u'_{p+1}$ と書く。第三の旗を探索する前に
\[
 W_p=(u'_p\otimes1)(\ker L)\subset V_{p+1},\qquad d_p=\dim W_p
\]
を求められる。$u'_p$ に $u_p$ の倍を加えても $W_p$ は変わらない。
$W_p$ の主元が「二番目の文字」で始まる先頭対を与えるので、
その文字の番号は $d_p+1$ に強制される。従って $d_p>2$ は直ちに排除できる。
さらに $W_p\subset\ker u_{p+1}$ であり、次の双対文字が $W_p$ 上で非零となることは、
次位置の第二文字を右に付けた列が主元となることと同値である。
次位置でも $d_{p+1}$ が定まるとき、
\[
         u'_{p+1}(W_p)\ne0 \quad\Longleftrightarrow\quad d_p>d_{p+1}
\]
が必要である。これは各辺の接方向と二次関係だけから四周期列を除外する検査となる。

\begin{corollary}\label{cor:pbw-no-double-cycle}
有限四辺の標準窓が周期的な左辞書順二次PBW表示を持つなら、
上の四周期列に沿う局所環の長さは少なくとも3である。
従って全ての局所因子が長さ1または2である有限辺を一つでも持つ場合もこの表示は存在しない
（他の三辺も有限とする）。
\end{corollary}
\begin{proof}
系\ref{cor:finite-edge-point-sequences}により、共有文字への隣接射影は同一の部分スキームへの同型である。
従って選んだ列に沿う四つの局所環は同型である。長さ1はすでに排除した。
長さ2なら各点は曲線状二重点であり、両射影は接線を同型に運ぶ。
辺 $p$ の接方向の両端には、非零スカラー倍を調整して
上の第二双対文字 $u'_p,u'_{p+1}$ を使える。

この辺の六関係の中で一次微分が消えるものは一次元であり、その生成元を $r_0$ とする。
$r_0(u'_p,u'_{p+1})$ は非零である。
実際、接方向を $\epsilon^3=0$ まで延ばすための方程式は、二次の補正を $b,b'$ として
\[
         b(D r)+b'(L r)=-r(u'_p,u'_{p+1})\qquad(r\in R_p)
\]
となる。左辺の像は $R_p^*$ の余次元1の空間であり、$r_0$ の消滅空間に一致する。
右辺が $r_0$ で零なら補正が存在する。しかし局所環 $k[t]/t^2$ から
$k[\epsilon]/\epsilon^3$ へ、$t$ の像の一次係数が非零となる写像は存在しない。
これは接方向の非零性に反する。$r_0\in\ker L$ なので $u'_{p+1}(W_p)\ne0$ である。
上の主元条件から四位置の全てで $d_p>d_{p+1}$ となり、周期性に反する。
\end{proof}

\paragraph{二重点を持つ正則窓への適用。}
標数零の代数閉体上、対称チャートで四辺の各々に被約点16個と長さ2の点2個を持つAS窓が存在する。
実際、四変数二次式の対称行列の空間 $\mathbb P^9$ で次の直線を取る：
\[
 M(s,t)=\begin{pmatrix}s I_2&t I_2\\ t I_2&0\end{pmatrix}.
\]
$t\ne0$ では可逆で、$P=[M(1,0)]$ だけが階数2である。
近傍で左上ブロックを $A$、右上を $B$、右下を $D$ とすれば、
階数2以下の軌跡の三つの局所式は $D-B^{\mathsf t}A^{-1}B=0$ である。
直線を含む一般の射影三平面を取り、追加する二方向を $N_1,N_2$ とする。
それらの右下対称ブロックは独立で、張る平面は $I_2$ を含まないようにできる。
局所変数を $t,u_1,u_2$ とすると三式の線型項は追加二方向だけに階数2を持つ。
その二成分を解けば $u_i=O(t^2)$ となり、残る成分の二次係数は
$-I_2$ の法線の商における像で非零である。従って $P$ の局所環の長さは2である。

直線を含む超平面系は直線外で基点自由である。その一般の六元で切ると、
Bertiniの定理により階数2の滑らかな六次元軌跡では $P$ 以外の交点は横断的であり、
直線と交わらない三次元の階数1軌跡は避ける。
また直線に属する二次式二本は共通因子を持たず、$\mathbb P^3$ での共通零点は曲線である。
追加の一般二次式二本はその曲線上に共通零点を持たない。
以上の非空開条件を同時に取れば、四式は定理\ref{thm:symmetric-frobenius-chart}の正則窓を与える。

辺スキームは二つの一次式の順序付き積がこの三平面に入る条件である。
積の射は階数2の上で次数2のエタール被覆である。実際、二因子は相異なり、
交換による商から階数2軌跡への写像は一意分解により有限双有理で、像は滑らかなので同型となる。
ゆえに $P$ の上に二つの二重点があり、他の全点は被約である。六つの双一次式の
有限完全交叉の長さは $\binom63=20$ なので被約点は16個である。
この例に系\ref{cor:pbw-no-double-cycle}を適用すると、辺が非被約でも
四周期の左辞書順二次PBW表示が存在しないことが分かる。

上の二方向を有理数で選んだ具体例は
\path{computations/verify_tangent_pencil_window.py}にある。四二次式の基点の不在を確かめ、
三次小行列式で切ったアフィン代数の辞書順基底から、分離座標の十次式とその導関数の最大公約元が
二重点を表す一次式となることを有理数上で検算する。無限遠に別の交点はない。

\begin{proposition}\label{prop:pbw-symmetric-square-flag}
標数零の代数閉体上で対称完全交叉チャートの四次元二次式系を $L$ とする。
全位置で共通の順序基底 $e_1<e_2<e_3<e_4$ を用い、番号もこの順序に合わせて
先頭対を $e_j e_i$ ($j>i$) にできるための必要十分条件は、次の二段の平方条件である。
$L$ がある非零な平方 $u_4^2$ を含み、$u_4=0$ に制限した三次元二次式系も非零な平方を含む。
\end{proposition}
\begin{proof}
$u_i$ を $e_i$ の双対変数とする。六関係は対称である。
指定した先頭対で被約行階段形を取ると、一行に他の下降対は現れず、
対称性により非対角項は指定対とその反対の和だけとなる。
加えられる平方は辞書順から $e_k^2$ ($k<j$) に限る。
従ってその直交する対称二次式系は、各 $k$ で
$u_k^2$ に $u_i u_j$ ($i<j$, $j>k$) の線型結合を加えた四式を基底とする。
逆にこの形なら、各下降対の対称和に下位の平方だけを加えて直交関係を作れる。
三角形の先頭語と三次元数20から前の系を適用できる。

この基底の形からまず $u_4^2$、次に制限後の $u_3^2$ を得る。
逆に、基点のない四式系に平方があれば、対応する超平面上への制限は基点のない三式系である。
実際、次元は高々3で、射影平面上の高々二つの二次式は共通零点を持つ。
次の平方を選んで同様に制限すると、二変数の基点のない鉛筆が残る。
その二元行列束の行列式は二次式で、代数閉体上では特異な非零要素を持ち、これも平方である。
最後の一変数に残るのは非零な平方だけである。逆順に持ち上げると
$q_k\equiv u_k^2\bmod(u_{k+1},\ldots,u_4)$、さらに後の基底式で高い番号の平方項を消せる。
これが求める形である。任意の位置別基底を許す場合にはこの必要性を主張しない。
\end{proof}

この表示と辺の有限性は両立する。同じ検算コードの
\texttt{triangular\_window} は上記三角形の有理四式を与える。
最後の式は $w^2$ である。先頭式の係数を1にしたチャートで
三次小行列式の辞書順基底の先頭項は $b,c,d^5$ となり、
最後の五次式は平方自由である。従って階数2以下の切断はこのチャートで被約五点である。
階数1なら全ての三次小行列式の微分が消えるから、この五点は全て階数2である。
残りのチャートのうち次またはその次の係数が先頭非零となるものは単位イデアルを与え、
残る支持は $w^2$ だけである。順序付きの辺には被約点十個と $(w,w)$ の一点があり、
有限完全交叉の長さから後者の局所長は10である。
